% Source: Corsin Nick/Class Notes/Week 8.pdf, pp. 5--9 (pp. 1--4 and 10--13 NOT yet read).
% ============================================================================
% PARTIAL TRANSCRIPTION. Pages 5-9 are transcribed below from the source. Pages
% 1-4 (priority page + start of the repetition quiz) and 10-13 (rest of the
% change-of-variables example, Feynman's trick) are NOT yet transcribed and are
% marked with TODO ainotes. See docs/07-todo.md.
% ============================================================================
\chapter{Week 8 --- Submanifold Quiz, Change of Variables \& Fubini}
\label{ch:week08}

\begin{ainote}[Chapter partially transcribed]
Pages 5--9 of the source are transcribed below. \textbf{Pages 1--4 and 10--13 are not yet
done} --- each gap is marked with its own TODO note at the point where it belongs.
\end{ainote}

\begin{ainote}[TODO --- quiz items 1--5 not transcribed]
The repetition quiz starts on p.~2; only items \textbf{6} and \textbf{7} (p.~5--6) are
transcribed here. One item in the untranscribed part is already known via \texttt{OQ-18}: a
bullet on p.~4 with $\det \Hess f(0) = 2-\alpha^2$ claims $\nabla f$ is a local diffeomorphism
\qt{for $\alpha \neq \sqrt2$}, but $2-\alpha^2$ vanishes at $-\sqrt2$ too, so it must read
$\alpha \neq \pm\sqrt2$.
\end{ainote}

% Source: Corsin Nick/Class Notes/Week 8.pdf, p. 5
\section{Repetition quiz}
\label{sec:week08_quiz}

\begin{question}[Quiz 6 --- graph over which variable?]
\label{q:week08_graph_over}
Let $V := \{(x,y) \in \mathbb{R}^2 : y^4+y^2 = x^3+x\}$. Decide:
\begin{enumerate}[label=\textbf{(\alph*)}]
  \item $V$ is a graph over $y$ in a neighbourhood of $(0,0)$.
  \item $V$ is a graph over $x$ in a neighbourhood of $(0,0)$.
\end{enumerate}
\end{question}

\begin{answer}
\textbf{(a) True.} Notice that $V = f^{-1}\{0\}$ for $f(x,y) := y^4+y^2-x^3-x$. Then
$$\Jac f(x,y) = (\partial_xf,\ \partial_yf) = (-3x^2-1,\ 4y^3+2y), \qquad \Jac f(0,0) = (-1,\ 0).$$
Since $\partial_xf(0,0) = -1 \neq 0$, the implicit function theorem
(\cref{thm:implicit_function_theorem}) applies with $x$ as the solved variable: in a
neighbourhood of $(0,0)$,
$$(x,y) \in V \iff f(x,y) = 0 \iff x = x(y).$$
The theorem does \textbf{not} apply for $y$, since $\partial_yf(0,0) = 0$.

\textbf{(b) False.} Solving the defining equation for $y^2$ by the quadratic formula gives
$$y^2 = -\tfrac12 + \sqrt{\tfrac14 + x^3 + x},$$
which has no solution for $x \in (-\varepsilon,0)$: there $x^3+x < 0$, so the square root falls
below $\tfrac12$ and the right-hand side is negative, while $y^2 \geq 0$. So $V$ has no points
at all just to the left of the origin, and cannot be a graph over $x$ there.
\end{answer}

% Sonnet 5 (Medium)
\begin{ainote}
This is \cref{thm:implicit_function_theorem} used as a decision procedure rather than an
existence theorem, and it is worth noticing that the two parts are settled by different means:
\textbf{(a)} by checking a derivative, \textbf{(b)} by an explicit computation showing the set is
empty on one side. The theorem gives no negative results --- failure of its hypothesis proves
nothing, exactly as discussed in \cref{sec:implicit_function_theorem}.
\end{ainote}

% Source: Corsin Nick/Class Notes/Week 8.pdf, p. 6
\begin{question}[Quiz 7 --- is an injective smooth curve a submanifold?]
\label{q:week08_injective_curve}
Let $\gamma : I \to \mathbb{R}^2$ be injective and smooth, $I$ an open interval. Is $\gamma(I)$
then a smooth submanifold of $\mathbb{R}^2$?
\end{question}

\begin{answer}
\textbf{False.} Consider the \newterm{figure eight} (\germanterm{Lemniskate})
$$\gamma : (0,2\pi) \to \mathbb{R}^2, \qquad t \mapsto (\sin t,\ \sin 2t).$$
This is injective and smooth, but there is no graphical representation of
$\operatorname{im}(\gamma)$ around $(0,0)$ --- the image crosses itself there.
\end{answer}

% Sonnet 5 (Medium) --- FIG-W08-01, the figure-eight from the source page.
\begin{center}
\begin{tikzpicture}[>=Latex, scale=1.5]
  \draw[->] (-1.35,0) -- (1.35,0) node[right, font=\scriptsize] {$x$};
  \draw[->] (0,-1.25) -- (0,1.25) node[above, font=\scriptsize] {$y$};
  \draw[blue!80!black, thick, variable=\t, domain=0:6.2832, samples=200, smooth]
    plot ({sin(\t r)}, {sin(2*\t r)});
  \draw[red!80!black, thick, dashed] (0,0) circle (0.34);
  \node[red!80!black, font=\scriptsize, right] at (0.42,0.30) {problem!};
  \draw[red!80!black, thin, ->] (0.42,0.26) -- (0.18,0.12);
\end{tikzpicture}
\end{center}

\begin{ainote}
The endpoints $t \to 0^+$ and $t \to 2\pi^-$ both send $\gamma(t) \to (0,0)$, so the origin is
approached from two different parameter ends without ever being attained --- which is how the
image can cross itself while $\gamma$ stays injective. The reason $\operatorname{im}(\gamma)$
fails to be a submanifold at the origin is exactly the component count of
\cref{ex:cross_not_submanifold}: deleting the crossing point leaves four pieces, a line minus a
point leaves two. Note also what this says about \cref{thm:local_parametrization} --- injectivity
and an injective differential are \emph{not} sufficient; the theorem also demands that the
parametrization be a \textbf{homeomorphism} onto its image, and that is what fails here.
\end{ainote}

% Sonnet 5 (Medium)
The rest of the week leaves differential topology behind and starts the theory of integration in
several variables.

% Source: Corsin Nick/Class Notes/Week 8.pdf, p. 7
\section{Multivariate integration}
\label{sec:multivariate_integration}

\subsection{Change of variables}
\label{sec:change_of_variables}

\begin{theorem}[Change of variables]
\label{thm:change_of_variables}
Suppose $U,V \subseteq \mathbb{R}^n$ are open sets and $\Phi : U \to V$ is a
$C^1$-diffeomorphism. If $A \subseteq U$ is Jordan measurable with $\overline{A} \subseteq U$,
and $f : A \to \mathbb{R}$ is continuous, then
$$\int_A f(x)\,dx = \int_{\Phi(A)} \frac{f\circ\Phi^{-1}(y)}{\big|\det \Jac\Phi\big(\Phi^{-1}(y)\big)\big|}\,dy.$$
\end{theorem}

\begin{ainote}[TODO --- Jordan measurability not transcribed]
\newterm{Jordan measurable} (\germanterm{Jordan-messbar}) is used here without being defined;
per \texttt{docs/03-topic-index.md} Week 8 covers Jordan measure (lecture ch.~13), presumably on
one of the pages not yet read. Definition still to be transcribed.
\end{ainote}

% Source: Corsin Nick/Class Notes/Week 8.pdf, p. 7
\subsection{Fubini's theorem}
\label{sec:fubini}

\begin{theorem}[Fubini]
\label{thm:fubini}
Let $X \subseteq \mathbb{R}^m$ and $Y \subseteq \mathbb{R}^n$ be intervals (\qt{boxes}), and let
$f : X\times Y \to \mathbb{R}$, $(x,y)\mapsto f(x,y)$, be continuous. Then
$$\int_{X\times Y} f(x,y)\,dx\,dy
  = \int_X\left(\int_Y f(x,y)\,dy\right)dx
  = \int_Y\left(\int_X f(x,y)\,dx\right)dy.$$
\end{theorem}

Here we look at the expression in brackets as a function, e.g.\
$x \mapsto \int_Y f(x,y)\,dy$.

% Source: Corsin Nick/Class Notes/Week 8.pdf, p. 8
\begin{example}[Fubini on the unit square]
\label{ex:fubini_unit_square}
$$
\begin{aligned}
\int_{[0,1]^2} x^{y+1}\,dx\,dy
  &= \int_0^1 dy \int_0^1 dx\ x^{y+1}
   = \int_0^1 dy\ \left[\frac{x^{y+2}}{y+2}\right]_0^1 \\
  &= \int_0^1 dy\ \frac{1}{y+2}
   = \log|y+2|\,\Big|_0^1
   = \log\!\left(\tfrac32\right).
\end{aligned}
$$
\end{example}

\begin{ainote}
Corsin writes the first line as $\int_0^1 dx \int_0^1 dy$, but the inner antiderivative taken on
the next line is in $x$ (it produces $x^{y+2}/(y+2)$, evaluated from $0$ to $1$). The
differentials are therefore swapped relative to the computation performed; the order is
corrected above and the result is unaffected. See \texttt{OQ-19}.
\end{ainote}

% Source: Corsin Nick/Class Notes/Week 8.pdf, pp. 8--9
\begin{example}[A domain that Fubini alone cannot handle]
\label{ex:change_of_variables_domain}
Consider
$$\int_A \frac{y}{\sqrt{x}}\,dy\,dx, \qquad
A := \left\{(x,y) \in \mathbb{R}^2 \ :\ x,y>0,\ \ 1 \leq \frac{y}{\sqrt x} \leq 2,\ \ 1 \leq xy \leq 2\right\}.$$
This is not easily integrated with Fubini alone. But we can apply a transformation of variables
to simplify the integration domain $A$: take
$$\Phi : U \to V, \qquad \begin{pmatrix}x\\y\end{pmatrix} \mapsto
  \begin{pmatrix}u(x,y)\\v(x,y)\end{pmatrix} = \begin{pmatrix}\tfrac{y}{\sqrt x}\\[2pt] xy\end{pmatrix},$$
with $U$, $V$ still to be determined --- they must be open.

\textbf{Check injectivity.} Suppose, for $x,y,x',y'>0$,
$$\textbf{(1)}\ \ \frac{y}{\sqrt x} = \frac{y'}{\sqrt{x'}},
\qquad\qquad \textbf{(2)}\ \ xy = x'y'.$$
From \textbf{(2)}, $\dfrac{x}{x'} = \dfrac{y'}{y}$. Inserting into \textbf{(1)},
$$\frac{1}{\sqrt x} = \frac{y'}{y}\cdot\frac{1}{\sqrt{x'}} = \frac{x}{x'}\cdot\frac{1}{\sqrt{x'}}
\quad\Longrightarrow\quad x^{3/2} = (x')^{3/2} \quad\Longrightarrow\quad x = x',$$
and consequently $\dfrac{y'}{y} = \dfrac{x}{x'} = 1$. So for $U := (0,\infty)^2$ and
$V := \Phi(U)$, the map $\Phi$ is a bijection between open sets.

We calculate the functional determinant $|\det \Jac\Phi|$. If it is strictly positive, we can
proceed with a change of variables, since $\Phi$ is then a local diffeomorphism and a bijection.
\end{example}

\begin{ainote}[TODO --- example unfinished]
The source continues on p.~10 with the determinant computation and the resulting integral over
the rectangle $[1,2]^2$. \textbf{Pages 10--13 are not yet transcribed}, so this example breaks
off mid-argument. Also still to come on those pages, per
\texttt{docs/03-topic-index.md}: \newterm{Feynman's trick} (differentiation under the integral
sign). One item there is already known via \texttt{OQ-17}: in the final integration on p.~13 the
$\alpha$-term reads $\tfrac{\pi}{2}\cdot\tfrac{\alpha}{1+\alpha^2}$ integrating to
$\tfrac{\pi}{4}\log(1+\alpha^2)$, both a factor $2$ off; corrected they are
$\tfrac{\pi}{4}\cdot\tfrac{\alpha}{1+\alpha^2}$ and $\tfrac{\pi}{8}\log(1+\alpha^2)$, which
makes his two bracket terms equal. His final answer $\tfrac{\pi\log 2}{8}$ is correct.
\end{ainote}

\exercisesheet{8}

\begin{ainote}[TODO --- recommended exercises deliberately skipped]
\textbf{Corsin's \textit{Recommended exercises} priority page for Week 8 (p.~1) has not been
transcribed}, and was skipped deliberately on this pass rather than overlooked. Every week from
4 onward carries one: a colour-coded table (blue \textbf{= important}, orange
\textbf{= semi-important}, red \textbf{= optional}) with a note per problem, transcribed into an
\texttt{ainote[Corsin's recommendations]} block --- see \cref{ch:week07} for the pattern.

Consequently \textbf{no exercises from sheet 8 are reproduced below}, since which ones to
reproduce is decided by that table. The sheet is \texttt{exercises/Ex8\_Analysis2\_eng.pdf}.
\end{ainote}

\section{Solutions}
\label{sec:week08_solutions}

\begin{ainote}[TODO --- empty pending transcription]
Empty until the exercise selection above is made.
\end{ainote}
